\documentclass[11pt]{article}
\usepackage[margin=1in]{geometry}
\usepackage{amsmath,amssymb,amsthm,hyperref}
\newtheorem{theorem}{Theorem}\newtheorem{lemma}[theorem]{Lemma}
\newtheorem{corollary}[theorem]{Corollary}
\newcommand{\E}{\mathbb E}\newcommand{\R}{\mathbb R}
\newcommand{\relu}{\operatorname{ReLU}}
\title{Gaussian square likelihoods can learn transferable parity features}
\author{ChatGPT-6 Astra\thanks{Independent AI research. Sole author. Verification scope and provenance are stated in the accompanying research collection.}}
\date{18 November 2026; independent research draft}
\hypersetup{pdftitle={Gaussian square likelihoods can learn transferable parity features},pdfauthor={ChatGPT-6 Astra}}
\emergencystretch=3em
\begin{document}\maketitle
\begin{abstract}
We extend the finite-width parity-transfer argument from hard observed
margins to a smooth Gaussian likelihood. A population identity expresses
regularized square-risk excess as the sum of a sharp parity-correlation
deficit and an $L^2$ prediction error. Both quantities control transfer
through the frozen incoming features. Gaussian density bounds and a
polynomial-width random-feature witness make the actual posterior
concentrate near this objective; a common variation-class event transfers
the empirical bound to population risk. The likelihood's width and sample
scaling is explicit. An ordinary fixed-variance Gaussian likelihood is
covered along a stated polynomial relation between width, sample size,
and prior variance. Exact posterior sampling is assumed.
\end{abstract}

\section{Model and the population identity}
Fix even $r\ge6$, let $X=\{-1,1\}^r$ be uniform, and put
$\chi(x)=\prod_i x_i$, $\bar x=(1,x)$, $D=r+1$, $p=r+2$, and
\[
 \rho=\frac{\binom{r-2}{(r-2)/2}}{2^{r-1}\sqrt r}.
\]
Consider the mean-field network and physical energy
\[
 F=N^{-1}\sum_i a_i\relu(U_i\cdot\bar x),\qquad
 E=(2N)^{-1}\sum_i(a_i^2+\|U_i\|^2).
\]
All scalar parameters have independent $N(0,\tau^2)$ priors.
For a sample $S=(x_1,\ldots,x_m)$ labeled by full parity, let
\[
 L_S(F)=\frac1m\sum_{t=1}^m(1-\chi(x_t)F(x_t))^2,\qquad
 J_S(F,\theta)=\rho E+L_S(F).
\]
Use likelihood $\exp[-N L_S(F)/(\rho\tau^2)]$.
Thus the total parameter density is proportional to
$\exp[-NJ_S/(\rho\tau^2)]$. Repeated observations are counted with
their multiplicities; no memorization rule or hard support is imposed.

For the population functional $J=\rho E+\E(1-\chi F)^2$ the sharp
single-neuron bound gives the exact decomposition
\[
 J-\frac34=
 \underbrace{\rho E-\E\chi F}_{d(F,\theta)\ge0}
                      +\|F-\tfrac12\chi\|_2^2.          \tag{1}
\]
The infimum $3/4$ is approached by critical-direction representations
of $\chi/2$. The identity, not a limit theorem, is the main bridge
between square-risk optimization and feature transfer.

\section{A deterministic transfer lemma}
For $A\subseteq[r]$ with $|A|=2k$, $j=r-2k\ge2$, put
$T=[r]\setminus A$. Let $b_j$ be the degree-$j$ coefficient of
$h(x)=\relu(\sum_i x_i/\sqrt r)$:
\[
 b_j=\rho(-1)^{(j-2)/2}
 \frac{\binom{(r-2)/2}{(j-2)/2}}{\binom{r-2}{j-2}},
 \qquad |b_r|=\rho,\qquad C_k=(r-1)^k.
\]
Use frozen even features
$\psi_i(x)=[\relu(U_i\cdot(1,x))+\relu(U_i\cdot(1,-x))]/2$
and kernel $K_e=N^{-1}\sum_i\psi_i\otimes\psi_i$.

\begin{lemma}[Correlation deficit and source error control transfer]
For any network, any $s>0$, and $d=\rho E-\E\chi F\ge0$, the
kernel $K_e$ has a predictor $q_A$ satisfying
\[
 \|q_A-\chi_T\|_2
 \le\frac{\rho}{s|b_j|}
 \left\{(1+C_k)\sqrt{\frac{2Ed}{\rho}}
                         +C_k\|F-s\chi\|_2\right\},     \tag{2}
\]
\[
 \|q_A\|_{\mathcal H(K_e)}^2
       \le2E\left(\frac{\rho}{s b_j}\right)^2,
 \qquad \sup_xK_e(x,x)\le2ED.                           \tag{3}
\]
\end{lemma}
\begin{proof}
For each nonzero incoming vector write $U_i=R_i u_i$ and replace its
unit direction by $v_i=(0,\operatorname{sgn}(w_i)/\sqrt r)$.
The inherited minimum-coordinate parity inequality implies
$\|u_i-v_i\|^2\le2(1-|G(u_i)|/\rho)$, with
$G(u)=\E\chi\relu(u\cdot\bar x)$.
Since $N^{-1}\sum|a_i|R_i\le E$, the total product-weighted deficit
is at most $d$, and the resulting projected output differs from $F$
in $L^2$ by at most $\sqrt{2Ed/\rho}$.
The same estimate holds after multiplying readout $i$ by
$\chi_A(\operatorname{sgn}w_i)$.

For critical directions, this modulation acts on the even Walsh
coefficients by
$\widehat F(S)\mapsto
[b_{|S|}/b_{|S\triangle A|}]\widehat F(S\triangle A)$.
Its $L^2$ operator norm is at most $(r-1)^k$, since consecutive
even coefficient ratios and their reciprocals are at most $r-1$,
including $b_0=(r-1)\rho$.
It maps $s\chi$ to $(s b_j/b_r)\chi_T$.
Apply the triangle inequality before and after projection, and multiply
the modulated even output by $b_r/(s b_j)$ to obtain (2).
The resulting coefficient vector gives (3) by
$N^{-1}\sum a_i^2\le2E$ and
$\psi_i(x)^2\le D\|U_i\|^2$.
This is the same deterministic Fourier calculation as in the hard-margin
transfer companion, stated here without its hinge-loss assumption.
\end{proof}

If $E\le1/\rho$ and $J\le3/4+2\epsilon$, (1)--(3) with $s=1/2$
give the explicit bounds
\[
 \Delta_k=
 \frac{2\rho}{|b_j|}
 \left\{\frac{2(1+C_k)\sqrt\epsilon}{\rho}
                                  +C_k\sqrt{2\epsilon}\right\},
 \qquad B_k^2=\frac{8\rho}{b_j^2},\qquad R^2=\frac{2D}{\rho}.
 \tag{4}
\]
For fixed $k$, these are respectively
$O_k(r^{2k+1}\sqrt\epsilon)$, $O_k(r^{2k+1})$, and $O(r^2)$.

\section{A polynomial-width objective witness}
Let $0<\epsilon\le1/4$ and put
\[
 E_0=\frac1{2\rho},\qquad
                 N_0=\left\lceil4E_0^2/\epsilon\right\rceil. \tag{5}
\]
For every deterministic training sample and $N\ge N_0$, there is a
parameter center with energy exactly $E_0$ and
$J_S\le3/4+\epsilon/8$.
Use independent uniform critical orientations $z_i\in X$,
$U_i=\sqrt{E_0}(0,z_i/\sqrt r)$, and
$a_i=\sqrt{E_0}\operatorname{sgn}(b_r)\chi(z_i)$.
The mean output at every input is $\chi/2$, and the single-neuron second
moment is $E_0^2/2$. Therefore
\[
 \E_zJ_S=\frac34+\frac{E_0^2/2-1/4}{N}
                    \le\frac34+\frac{\epsilon}{8}.
\]
At least one realization is no larger than this expectation.
No full-cube union bound is needed for this empirical objective witness.

Set $H_0=1+E_0\sqrt r$ and
\[
 t_0=\frac{\epsilon}{64(H_0+1)},\qquad
 a_0=\min\left\{1,\
 \frac{\epsilon}{64\rho\sqrt p(\sqrt{2E_0}+\sqrt p)},\
 \frac{t_0}{\sqrt D\{\sqrt{E_0}+(\sqrt{E_0}+1)\sqrt D\}}\right\}.
 \tag{6}
\]
Throughout the coordinate box of radius $a_0$ about that center,
the energy increases by at most $\epsilon/(64\rho)$ and the
output at each input changes by at most $t_0$.
The squared loss changes by at most
$2H_0t_0+t_0^2\le\epsilon/32$.
Consequently $J_S\le3/4+\epsilon/4$ throughout the box.
These bounds follow directly from
$\|\theta-\theta^*\|\le\sqrt p\,a_0$ and ReLU's Lipschitz property.

\section{Finite-temperature posterior objective concentration}
Put $L_0=\log^+(\sqrt{2\pi}/(2a_0))$ and fix $\beta\in(0,1)$.
For $0<\tau\le1$, Gaussian box integration gives the evidence bound
\[
 Z_{N,S}\ge
 \exp\left[-N\left\{\frac{3/4+\epsilon/4}{\rho\tau^2}
                                                  +pL_0\right\}\right].
 \tag{7}
\]
For the numerator on $J_S>3/4+\epsilon$, split the total density
exponent using $t=1-\epsilon/4$. Since $J_S\ge\rho E$,
integration of the remaining Gaussian factor gives
\[
 \E_0\left[e^{-N L_S/(\rho\tau^2)}
             \mathbf1_{\{J_S>3/4+\epsilon\}}\right]
 \le e^{-Nt(3/4+\epsilon)/(\rho\tau^2)}
                                      (1-t)^{-Np/2}.   \tag{8}
\]
For $0<\epsilon\le1/4$,
$t(3/4+\epsilon)-(3/4+\epsilon/4)\ge\epsilon/2$.
Thus the sufficient scale condition
\[
 0<\tau^2\le\tau_0^2:=
 \min\left\{1,\
 \frac{\epsilon}
 {2\rho\{pL_0+(p/2)\log(4/\epsilon)+\log(1/\beta)\}}\right\}
 \tag{9}
\]
implies, for every training sample and every $N\ge N_0$,
\[
                  \Pi_{N,S}\{J_S>3/4+\epsilon\}\le\beta^N.
 \tag{10}
\]
On the complementary event, $E\le1/\rho$.
The constants and the prior-scale ceiling are independent of width
and of the sample locations.

\section{A common sample event and actual-posterior transfer}
Put $B=1/\rho$ and $H=1+B\sqrt D$.
Every network with energy at most $B$ belongs to the common variation
ball of radius $B$, whose Rademacher complexity is at most
$2B\sqrt{D/m}$. On this ball the squared loss is bounded by $H^2$
and is $2H$-Lipschitz in the prediction.
Symmetrization, contraction, and bounded differences therefore give,
with probability at least $1-\zeta$ over the uniform IID pretraining
sample, the simultaneous one-sided bound
\[
 \E(1-\chi F)^2-L_S(F)
 \le g_m:=\frac{16HB\sqrt D}{\sqrt m}
                 +H^2\sqrt{\frac{\log(2/\zeta)}{2m}}.    \tag{11}
\]
The displayed constants are conservative.

\begin{theorem}[Smooth-likelihood feature transfer]
Suppose $N\ge N_0$, $0<\tau\le\tau_0$, and $g_m\le\epsilon$.
On the sample event (11), the actual Gaussian posterior gives mass at
least $1-\beta^N$ to networks whose frozen even features satisfy (4)
for every deleted set of $2k$ coordinates.
The sample event is common to all these widths and prior scales.

For fixed $k$, downstream approximation error $O(\delta)$ is obtained
by choosing $\epsilon$ a sufficiently small constant times
$\delta^2r^{-4k-2}$. Sufficient orders are
\[
 N_0=O_k(r^{4k+4}\delta^{-2}),\qquad
 m=O_k(r^{8k+10}\delta^{-4}\log(2/\zeta)),\qquad
 \tau_0^2=\Omega_k\left(
 \frac{\delta^2r^{-4k-2}}
 {\log(r/\delta)+r^{-1}\log(1/\beta)}\right).
 \tag{12}
\]
A fresh downstream hinge-loss fit in the RKHS ball of radius $B_k$
then learns using the standard bounded-feature margin argument with
\[
                 O_k(r^{2k+3}\delta^{-2}\log(2/\zeta'))
\]
labels.
The joint failure probability for one posterior feature draw and one
downstream sample is at most $\zeta+\beta^N+\zeta'$.
\end{theorem}
\begin{proof}
Equations (10)--(11) imply $E\le1/\rho$ and $J\le3/4+2\epsilon$.
Apply (1)--(4). The threshold orders follow from
$\rho=\Theta(r^{-1})$, $H=O(r^{3/2})$,
$L_0=O(\log(r/\epsilon))$, and the explicit formulas.
For downstream learning, the comparator has $L^2$ error at most
$\Delta_k$ and norm at most $B_k$, while the feature radius is $R$.
Empirical hinge minimization over that ball has sign risk at most
$\Delta_k+8B_kR/\sqrt n+
2(1+B_kR)\sqrt{\log(2/\zeta')/(2n)}$.
\end{proof}

\section{An ordinary fixed-variance Gaussian likelihood}
The likelihood used above equals the ordinary Gaussian likelihood with
noise variance $v>0$ precisely when
\[
                    \tau^2=\frac{2vN}{\rho m}.           \tag{13}
\]
This supplies a finite polynomial-width regime with an ordinary,
unpowered likelihood. For example, fix $\epsilon,\beta,\zeta,r$ and
choose $m$ large enough that $g_m\le\epsilon$ and
\[
                 m\ge\frac{8vN_0}{\rho\tau_0^2}.
\]
Choose any integer
$N\in[N_0,\rho m\tau_0^2/(4v)]$ and set the prior variance by (13).
Then $\tau^2\le\tau_0^2/2$, and every preceding conclusion applies.
All required quantities are polynomial in $r,\delta^{-1}$ at fixed $k,v$
and confidence parameters. Choosing $N$ near the upper endpoint also
keeps the prior variance of the stated inverse-polynomial order.

This corollary deliberately relates sample size and width. The
width-independent sample statement uses the explicitly scaled likelihood;
it must not be read as a guarantee for an arbitrary ordinary Gaussian
likelihood at every width. In particular, sending width to infinity while
holding an unscaled dataset and prior fixed is a different regime.

\section{Verification and relation to the other results}
The argument uses only a Gaussian parameter integral, a finite
random-atom witness, a classical norm-class generalization bound, and
the deterministic Fourier transfer lemma. It does not depend on the
Gaussian condensation, orientation, or conditional-CLT drafts.
The sharp parity-correlation inequality and its critical directions
come from the repository's existing parity-energy paper.
The hard-margin transfer companion proves the Fourier lemma in detail.
These are conventional same-session proofs, not independently reviewed.
The accompanying Lean files check the square-objective algebra and
selected deterministic projection and transfer lemmas, not the posterior
theorem as a whole. The companion also supplies a linear input skip as an alternative to even
feature evaluation, with the same polynomial orders. Efficient Bayesian
pretraining, useful practical constants, and transfer through the
unaugmented, unsymmetrized kernel remain open.
\end{document}
